\documentclass[10pt,a4paper]{amsart}
\usepackage{amsmath,amssymb,amsthm,graphicx}
% Retained minimal preamble for the Stacks Project chapter more-morphisms.tex (More on Morphisms).
% Verbatim excerpt of the chapter preamble: Xy-pic setup (native preamble lines 14-18) and the native
% theorem declarations with the mathematical macros (native preamble lines 154-217). Copied mechanically;
% nothing added, removed or redefined. The wrapper class and the amsmath/amssymb/amsthm/graphicx base
% packages are supplied by the assembler, so the class/externaldocument part of the native preamble is omitted.
% For commutative diagrams we use Xy-pic
\usepackage[all]{xy}

% We use 2cell for 2-commutative diagrams.
\xyoption{2cell}
% Theorem environments.
%
\theoremstyle{plain}
\newtheorem{theorem}[subsection]{Theorem}
\newtheorem{proposition}[subsection]{Proposition}
\newtheorem{lemma}[subsection]{Lemma}

\theoremstyle{definition}
\newtheorem{definition}[subsection]{Definition}
\newtheorem{example}[subsection]{Example}
\newtheorem{exercise}[subsection]{Exercise}
\newtheorem{situation}[subsection]{Situation}

\theoremstyle{remark}
\newtheorem{remark}[subsection]{Remark}
\newtheorem{remarks}[subsection]{Remarks}

\numberwithin{equation}{subsection}

% Macros
%
\def\lim{\mathop{\mathrm{lim}}\nolimits}
\def\colim{\mathop{\mathrm{colim}}\nolimits}
\def\Spec{\mathop{\mathrm{Spec}}}
\def\Hom{\mathop{\mathrm{Hom}}\nolimits}
\def\Ext{\mathop{\mathrm{Ext}}\nolimits}
\def\SheafHom{\mathop{\mathcal{H}\!\mathit{om}}\nolimits}
\def\SheafExt{\mathop{\mathcal{E}\!\mathit{xt}}\nolimits}
\def\Sch{\mathit{Sch}}
\def\Mor{\mathop{\mathrm{Mor}}\nolimits}
\def\Ob{\mathop{\mathrm{Ob}}\nolimits}
\def\Sh{\mathop{\mathit{Sh}}\nolimits}
\def\NL{\mathop{N\!L}\nolimits}
\def\CH{\mathop{\mathrm{CH}}\nolimits}
\def\proetale{{pro\text{-}\acute{e}tale}}
\def\etale{{\acute{e}tale}}
\def\QCoh{\mathit{QCoh}}
\def\Ker{\mathop{\mathrm{Ker}}}
\def\Im{\mathop{\mathrm{Im}}}
\def\Coker{\mathop{\mathrm{Coker}}}
\def\Coim{\mathop{\mathrm{Coim}}}

% Boxtimes
%
\DeclareMathSymbol{\boxtimes}{\mathbin}{AMSa}{"02}

%
% Macros for moduli stacks/spaces
%
\def\QCohstack{\mathcal{QC}\!\mathit{oh}}
\def\Cohstack{\mathcal{C}\!\mathit{oh}}
\def\Spacesstack{\mathcal{S}\!\mathit{paces}}
\def\Quotfunctor{\mathrm{Quot}}
\def\Hilbfunctor{\mathrm{Hilb}}
\def\Curvesstack{\mathcal{C}\!\mathit{urves}}
\def\Polarizedstack{\mathcal{P}\!\mathit{olarized}}
\def\Complexesstack{\mathcal{C}\!\mathit{omplexes}}
% \Pic is the operator that assigns to X its picard group, usage \Pic(X)
% \Picardstack_{X/B} denotes the Picard stack of X over B
% \Picardfunctor_{X/B} denotes the Picard functor of X over B
\def\Pic{\mathop{\mathrm{Pic}}\nolimits}
\def\Picardstack{\mathcal{P}\!\mathit{ic}}
\def\Picardfunctor{\mathrm{Pic}}
\def\Deformationcategory{\mathcal{D}\!\mathit{ef}}
\title{Universal openness is tested on affine-space extensions and on locally finitely presented base changes}
\author{The Stacks Project Authors}
\date{stacks-project:more-morphisms.tex:local-snapshot-20261011; GNU FDL 1.2 or later}
\begin{document}
\maketitle
\noindent\textit{Editorial note.} Statement span is the whole lemma environment lines 21725-21737 with the three enumerated conditions; proof span is lines 21739-21785.\par\medskip
\noindent\textit{Editorial note.} This card and S3-morph-char-univ-open-finite are deliberately both counted: the two lemmas are distinct statements in the same section, and the second one's proof reuses the polynomial-square trick rather than restating the first.\par\medskip
\noindent\textit{Editorial note.} The proof contains one xymatrix display; the retained preamble provides the Xy-pic setup.\par\medskip
\noindent\textit{Editorial note.} The proof cites no label and no bibliography entry, so no external\_references are needed.\par\medskip
\noindent\textit{Editorial note.} extra\_packages is empty on purpose: the retained preamble already loads xy with options and a second declaration clashes.\par\medskip
\noindent\textit{Allowed standard prerequisites (editorial).}\begin{itemize}
\item Explicit prerequisite boundary. Definition of a universally open morphism, base change of schemes, the principal opens D(g) as a basis, and the affine-space extension A\textasciicircum{}n x\_S X. The two reductions are retained in full.
\end{itemize}\medskip
\begin{lemma}
\label{lemma-test-universally-open}
Let $f : X \to S$ be a morphism of schemes.
The following are equivalent
\begin{enumerate}
\item $f$ is universally open,
\item for every morphism $S' \to S$ which is locally of finite presentation
the base change $X_{S'} \to S'$ is open, and
\item for every $n$ the morphism
$\mathbf{A}^n \times X \to \mathbf{A}^n \times S$
is open.
\end{enumerate}
\end{lemma}

\begin{proof}
It is clear that (1) implies (2) and (2) implies (3).
Let us prove that (3) implies (1).
Suppose that the base change $X_T \to T$ is not open
for some morphism of schemes $g : T \to S$.
Then we can find some affine opens
$V \subset S$, $U \subset X$, $W \subset T$
with $f(U) \subset V$ and $g(W) \subset V$
such that $U \times_V W \to W$ is not open.
If we can show that this implies
$\mathbf{A}^n \times U \to \mathbf{A}^n \times V$
is not open, then $\mathbf{A}^n \times X \to \mathbf{A}^n \times S$
is not open and the proof is complete. This reduces us
to the result proved in the next paragraph.

\medskip\noindent
Let $A \to B$ be a ring map such that $A' \to B' = A' \otimes_A B$
does not induce an open map of spectra for some $A$-algebra $A'$.
As the principal opens give a basis for the topology of $\Spec(B')$
we conclude that the image of $D(g)$ in $\Spec(A')$
is not open for some $g \in B'$. Write
$g = \sum_{i = 1, \ldots, n} a'_i \otimes b_i$
for some $n$, $a'_i \in A'$, and $b_i \in B$.
Consider the element $h = \sum_{i = 1, \ldots, n} x_i b_i$
in $B[x_1, \ldots, x_n]$. Assume that $D(h)$ maps to
an open subset under the morphism
$$
\Spec(B[x_1, \ldots, x_n]) \longrightarrow \Spec(A[x_1, \ldots, x_n])
$$
in order to get a contradiction. Then $D(h)$ would map surjectively
onto a quasi-compact open $U \subset \Spec(A[x_1, \ldots, x_n])$.
Let $A[x_1, \ldots, x_n] \to A'$ be the $A$-algebra homomorphism
sending $x_i$ to $a'_i$. This also induces a $B$-algebra
homomorphism $B[x_1, \ldots, x_n] \to B'$ sending $h$ to $g$.
Since
$$
\xymatrix{
\Spec(B[x_1, \ldots, x_n]) \ar[d] &
\Spec(B') \ar[l] \ar[d] \\
\Spec(A[x_1, \ldots, x_n]) &
\Spec(A') \ar[l]
}
$$
is cartesian the image of $D(g)$ in $\Spec(A')$ is equal to the
inverse image of $U$ in $\Spec(A')$ and hence open which is
the desired contradiction.
\end{proof}

\end{document}
